\title{Byte Latent Transformer: Patches Scale Better Than Tokens}

\begin{abstract}
We present the Byte Latent Transformer (\textsc{BLT}), a novel architecture for large language models operating directly at the byte level. For the first time, \textsc{BLT} achieves parity with traditional tokenization-based LLMs in performance at scale, while offering notable gains in inference speed and resilience. \textsc{BLT} processes byte sequences by grouping them into adaptively sized patches that constitute the fundamental computational elements of the model. These patches are determined according to the entropy of upcoming bytes, thereby concentrating computational effort and model capacity where the complexity of the input is greatest. We report the first study to investigate scaling of byte-level models with controlled computational cost, extending up to 8B parameters and 4T training bytes. Our experiments confirm that scaling models on unprocessed byte data—eschewing static vocabularies—is practical and effective. The approach also leads to improved efficiency during both training and inference, as the model forms longer patches in more predictable contexts, coupled with observable benefits in reasoning abilities and robust generalization to rare cases. In summary, under equivalent inference budgets, \textsc{BLT} delivers substantially enhanced scaling characteristics compared to models relying on tokenization, by expanding both patch and model size simultaneously.
\end{abstract}

\section{Introduction}
\label{section:intro}

We present the Byte Latent Transformer~(\textbf{BLT}), a novel tokenizer-free model that processes data directly at the byte level. For the first time, BLT attains parity in performance with conventional models that rely on tokenization, while offering notable gains in computational efficiency and resilience (\S\ref{sec:robustness}). In prevailing large language models (llms), training is nearly fully end-to-end except for the tokenization stage, where bytes are heuristically aggregated into a predetermined vocabulary of tokens. This tokenization process can introduce biases in text representation, resulting in limitations such as sensitivity to domain and modality~\citep{dagangetting}, heightened vulnerability to input perturbations~(\S\ref{sec:robustness}), insufficient encoding of orthographic structures~\citep{edman2024cute}, and disparities across languages~\citep{liang2023xlm,petrov2024language,limisiewicz2024myte}.

Tokenization has historically played a critical role since training llms directly on byte sequences incurs substantial computational expense, primarily because of increased sequence lengths~\citep{xue2022byt5}. To address this, earlier research has leveraged more computationally efficient self-attention mechanisms~\citep{el2020characterbert, clark2022canine} or adopted attention-free model architectures~\citep{wang2024mambabyte}~(\S\ref{section:related_work}). Nevertheless, these approaches mainly reduce costs when training \textit{small models}. For large-scale models, most of the computational burden arises from processing each byte through expansive feed-forward network layers, with the attention component contributing only a minor portion to the overall cost.

To optimize the use of computational resources, we introduce a dynamic and trainable technique for aggregating bytes into \textit{patches}~(\S\ref{section:patching}), accompanied by a novel model design that integrates both byte-level and patch-level information. In contrast to traditional tokenization, {\textsc BLT} does not utilize a predefined vocabulary for constructing patches. Instead, any selected grouping of bytes can be transformed into latent patch embeddings through efficient, parameterized encoder and decoder components. Our findings demonstrate that this strategy enables a \textit{more} effective distribution of computation compared to approaches based on tokenization.

Conventional tokenization-based llms devote an equal amount of computational resources to every token. This approach prioritizes simplicity at the cost of efficiency, as the use of compression heuristics to define tokens does not always align with the actual complexity involved in predicting those tokens. Our proposed architecture is grounded in the principle that computational effort should be directed adaptively, focusing more on the areas where it is most necessary. For instance, extensive transformer computation is usually redundant for generating the final characters of many words, which typically present straightforward, low-entropy decisions, especially when contrasted with the complexity of initiating a new sentence. This principle is embodied in the design of {\textsc BLT}~(\S\ref{section:architecture}), which incorporates three distinct transformer modules: two compact, byte-level \textit{local models}, and a more substantial, global \textit{latent transformer}~(\autoref{fig:arch_high_level}). 

To facilitate the dynamic assignment of computation, including how bytes are batched into patches, {\textsc BLT} partitions data by evaluating the entropy associated with predicting the next byte. This process yields contextual clusters of bytes, each characterized by a similar level of informational content.

We conduct the first scaling analysis of byte-level models regulated by flop count, extending up to 8B parameters and processing 4T bytes during training. This demonstrates the feasibility of training large-scale, end-to-end models directly from bytes, thus eliminating the need for fixed-vocabulary tokenization. In summary, {\textsc BLT} achieves training flop-controlled performance on par with Llama 3 while reducing inference computational cost by up to 50\%~(\S\ref{section:scaling}).\footnote{The computational demand is quantified by tallying the number of Floating Point OPerations (flops) required.} Additionally, our findings indicate that utilizing raw byte inputs markedly enhances performance in modeling infrequent or atypical data examples. Compared to models that rely on tokenization, {\textsc BLT} is substantially more robust to input noise and demonstrates superior character-level reasoning, which we illustrate through tasks involving orthographic knowledge, phonology, and machine translation for low-resource languages~(\S\ref{sec:robustness}).

Moreover, with the {\textsc BLT} architecture, both model size and input patch size can be increased concurrently while keeping the inference flop requirements constant. Employing larger patch sizes generally decreases computational expenditure, thereby allowing the reallocation of these savings towards expanding the global latent transformer component, which in turn is invoked less frequently. Our inference-flop constrained scaling studies~(\autoref{fig:fixed_inference_scaling}) reveal substantially more favorable scaling properties compared to tokenization-based approaches.

To conclude, the principal contributions of this work are as follows: 1) We present {\textsc BLT}, a byte latent llm architecture that enables adaptive computation allocation to enhance flop efficiency, 2) We demonstrate that our approach matches the training flop-to-performance ratio of Llama 3 up to 8B parameters, with the flexibility to accept slight decreases in evaluation scores in exchange for up to 50\% improvements in flop efficiency, 3) {\textsc BLT} models offer a novel scaling paradigm for llms, facilitating increases in model capacity without exceeding a fixed inference cost, 4) Our results indicate that {\textsc BLT} models possess greater resilience to noisy inputs and can capture sub-word information often overlooked by token-based llms. The training and inference code for {\textsc BLT} is available at~\url{https://github.com/facebookresearch/blt}.

\section{Patching: From Individual Bytes to Groups of Bytes}
\label{section:patching}

Dividing bytes into \textit{patches} enables {\textsc BLT} to adjust computational resources in response to contextual needs. As illustrated in Figure~\ref{fig:patching_types}, there are multiple strategies for partitioning bytes into patches. More formally, a patching function $f_p$ takes an input byte sequence $\pmb{x}=\{x_i,|i=1,\ldots n\}$ of length $n$ and segments it into a series of $m < n$ patches, represented as $\pmb{p}=\{p_j|j=1,\ldots,m\}$, by associating each $x_i$ with an indicator from \{0,1\}; a value of 1 signifies the initiation of a new patch. In both token-based and patch-based architectures, the overall compute required primarily hinges on the number of iterations the main Transformer module performs. For {\textsc BLT}, this corresponds to the total patches produced when encoding the input under a particular patching method. Thus, the mean patch length—referred to as \textit{patch size}—is central to estimating the computational cost for training and inference with any given patching strategy~(\S\ref{section:flops}).

We proceed by introducing three specific patching functions: segmenting with a constant byte count per patch~(\S\ref{section:static-patch}), applying whitespace-based patching~(\S\ref{section:white-patch}), and using adaptive patching informed by entropy values from a compact byte-level language model~(\S\ref{section:dyn-patch}). Lastly, we address approaches for incremental patching and clarify the distinctions between tokenization and patching~(\S\ref{section:bpe-patch}).

\subsection{Strided Patching Every K Bytes}
\label{section:static-patch}

A common and intuitive method for organizing bytes is to divide them into uniformly sized patches of length $k$, as implemented in MegaByte~\citep{yu2023megabyte}. Utilizing a fixed stride is both practical and straightforward for both training and inference, allowing for easy adjustment of the average patch size and thereby offering convenient control over computational cost. Despite these advantages, fixed-size patching has notable limitations. Most significantly, computational resources are distributed uniformly rather than in proportion to content complexity: for example, one might squander a transformer step $j$ merely predicting whitespace in code, or conversely, insufficiently allocate computation to information-rich segments, such as mathematical content. Furthermore, this method often results in inconsistent and context-insensitive partitioning of analogous byte sequences, with instances like the same word potentially being segmented in various ways.

\subsection{Space Patching}
\label{section:white-patch}

\citet{slagle2024spacebyte} introduces a straightforward yet effective enhancement to strided patching by initiating new patches immediately after any space-like byte\footnote{
    Space-like bytes refer to any byte not classified as a Latin character, a digit, or a utf-8 continuation byte. Furthermore, each patch must include at least one byte that is not space-like. }. These space-like bytes frequently correspond to linguistically meaningful boundaries in numerous languages. In the approach called Space patching, the latent transformer dedicates an additional computational step (i.e., increased flops) to every word encountered. This design guarantees that words are consistently patched throughout various sequences, while computational resources are focused on more challenging predictions, which typically occur after space-like characters. To illustrate, generating the initial byte in the answer to the query ``Who composed the Magic Flute?\uline{\hspace{.6em}}'' poses greater difficulty compared to the subsequent bytes after ``M'', since this first letter sharply limits the set of plausible completions—making ``Mozart'' relatively trivial to predict once the initial character is known. Nevertheless, the space patching strategy exhibits limitations: it does not adapt seamlessly to all languages or application areas, and crucially, it lacks the flexibility to alter patch size. In the following section, we propose an alternative patching technique based on the observation that the initial bytes of words are generally the most challenging to model, and that also provides a flexible way to adjust patch sizes.

\subsection{Entropy Patching: Using Next-Byte Entropies from a Small Byte LM}
\label{section:dyn-patch}

Instead of depending on heuristic rules like whitespace to determine boundaries, our method adopts a data-centric strategy to pinpoint where next-byte predictions are most uncertain. We propose a technique called \textit{entropy patching}, which leverages entropy measurements to determine the patch boundaries.

To achieve this, we train a compact, byte-level auto-regressive language model using the {\textsc BLT} training corpus, and for each position, we calculate the next byte entropies according to the LM distribution $p_{e}$ across the byte vocabulary $\mathcal{V}$:

\begin{equation}
H(x_i) = \sum_{v \in \mathcal{V}} p_{e}(x_i=v|\pmb{x}_{<i}) \log p_{e}(x_i=v|\pmb{x}_{<i})
\end{equation}

We investigate two approaches for determining patch boundaries using entropies $H(x_i)$. The first method selects points whose entropy values surpass a fixed global threshold, as depicted in~\autoref{fig:patching}. The second method locates points where the entropy is significantly greater compared to the preceding value. This latter strategy can also be viewed as detecting locations where the otherwise roughly monotonically decreasing entropy pattern within a patch is disrupted.

\begin{align*}
    \text{Global Constraint}\hspace{1cm}H(x_t)& > \theta_g\\
    \text{Approx. Monotonic Constraint}\hspace{1cm}H(x_t)& - H(x_{t-1}) > \theta_r
\end{align*}

Patch boundaries are determined in a simple preprocessing phase that takes place as data is loaded. This contrasts with the approach in~\citet{nawrot-etal-2023-efficient}, where a classifier is trained specifically to detect patch boundaries using entropy-based predictions. In our experiments~(\S\ref{section:experiments}), we evaluate and compare these two strategies for differentiating between bytes of low and high entropy.

\subsection{The Byte-Pair Encoding (BPE) Tokenizer and Incremental Patching}
\label{section:incremental-patching}
\label{section:bpe-patch}

A number of contemporary llms, such as our baseline Llama 3, rely on a subword tokenizer, notably bpe~\citep{gage1994new, sennrich-etal-2016-neural}. Throughout this work, we define ``tokens'' as groups of bytes selected from a \textit{finite} vocabulary that is established before model training, contrasting them with ``patches,'' which are sequences dynamically formed without reliance on a predefined vocabulary. One important distinction is that, when using tokens, the model does not interact directly with the raw byte-level features.

An important advancement offered by {\textsc BLT} relative to tokenization-based architectures is its novel redefinition of the balance between vocabulary size and computational requirements. In conventional large language models, expanding the vocabulary leads to a greater average token length, thereby reducing the number of inference steps. However, this also results in an expanded output dimensionality for the model’s final projection layer. Consequently, tokenization-based systems are constrained in their ability to substantially alter both token granularity and inference complexity. As an illustration, Llama 3 raises the average token length from 3.7 to 4.4 bytes, but this enhancement necessitates a fourfold enlargement of the embedding table when compared to Llama 2.

During generation, {\textsc BLT} must determine at each point in the byte sequence whether it has reached a patch boundary, as this choice dictates whether the Latent Transformer should be applied for further computation. Crucially, this boundary detection must be performed without knowledge of any future bytes, since the remainder of the sequence has not yet been produced. Consequently, the patching method must not rely on lookahead information to divide the byte sequence. More precisely, a patching scheme $f_p$ is said to exhibit incremental patching if the following holds:
$$f_p(\pmb{x}_{<i}) = f_p(\pmb{x})_{<i}$$
bpe does not fulfill the criteria for incremental patching, since an identical prefix could result in varying tokenizations based on its subsequent context, thereby violating the requirement above\footnote{It is possible to adapt bpe into an incremental patching scheme by introducing a dedicated delimiter token to denote patch boundaries, though this would extend the overall length of the byte sequence.}.

\section{{\textsc BLT} Architecture}
\label{section:architecture}
The {\textsc BLT} framework consists of a high-capacity, global autoregressive language model responsible for processing patch representations, augmented by two smaller local models: one for converting byte sequences into patches, and another for reconstructing byte sequences from patch representations~(\autoref{fig:arch_high_level}).

\subsection{Latent Global Transformer Model}

\textit{The Latent Global Transformer} refers to an autoregressive transformer architecture, denoted as $\mathcal{G}$, comprising $l_{\mathcal{G}}$ layers. This model translates a sequence of latent input patch embeddings, $p_j$, into a corresponding sequence of output patch embeddings, $o_j$. In this work, patch indices are represented by the subscript $j$, while byte indices are denoted with $i$. The global transformer employs a block-causal attention mask~\citep{dubey2024llama}, ensuring that attention is limited to the current and preceding patches within the same document. Since the global model dominates computational costs in both pre-training and inference stages, modulating when the model is utilized provides a mechanism to allocate and adapt computation dynamically to different segments of the input or output, contingent on their respective complexities.

\subsection{Local Encoder}
\label{section:local}

\textit{The Local Encoder Model}, represented as $\mathcal{E}$, is a compact transformer-style network comprising $l_{\mathcal{E}}$ layers, with $l_{\mathcal{E}}$ significantly smaller than $l_{\mathcal{G}}$. Its principal function is to rapidly transform an input byte sequence, $b_i$, into rich patch-level representations, denoted $p_j$. Unlike standard transformer designs, this model interleaves a cross-attention mechanism after each transformer layer. The cross-attention layers are responsible for aggregating byte-level encodings into patch representations, as shown in~(\autoref{fig:crossattn}).

Initially, each byte $b_i$ is converted to a vector using an embedding matrix in $\mathbb{R}^{256 \times h_{\mathcal{E}}}$, resulting in embeddings $x_i$. Optionally, these vectors may be further supplemented with hash-embeddings, as explained in~(\S\ref{sec:hash-emb}).

Subsequently, multiple blocks, each consisting of a transformer layer followed by cross-attention, iteratively refine the representations, ultimately yielding patch representations $p_i$. These $p_i$ vectors are then consumed by the global transformer, $\mathcal{G}$.

The transformer component of $\mathcal{E}$ employs a \textit{local block causal} attention pattern: each byte is allowed to attend only to a window of $w_{\mathcal{E}}$ preceding bytes. This window may span across changing patch boundaries but is constrained not to extend over document boundaries. The sections that follow elaborate on both the embedding process and the architecture of the cross-attention component.

\subsubsection{Encoder Hash n-gram Embeddings}
\label{sec:ngram-emb}
\label{sec:hash-emb}
An essential aspect of generating rich and resilient representations at every position $i$ involves embedding context from prior bytes. Within {\textsc BLT}, this is realized by representing not only the single byte $b_i$, but also its occurrence within a sequence of bytes—namely, a byte n-gram. To this end, at each step $i$, we begin by generating byte-grams

\begin{align}
    g_{i,n}=\{b_{i-n + 1},\ldots, b_i\}
\end{align}

for each byte position $i$ and for $n$ ranging from three to eight.\footnote{
We exclude byte-grams of length $n$ or greater if $i<n$. }

Next, we propose hash-based $n$-gram embeddings, in which all byte $n$-grams are mapped to a position within a fixed-size embedding matrix $E_{n}^{hash}$ using a hash function, for every $n \in\{3,4,5,6,7,8\}$~\citep{bai2010learning}. The embedding derived in this way is subsequently combined with the individual byte embedding, after which the sum is normalized and used as input for the local encoder. The enriched embedding is computed as follows:

\begin{align}
e_i &= x_i + \sum_{n = 3,...,8}E_{n}^{hash}(\text{Hash}(g_{i,n})) \\
\text{where, Hash}(g_{i,n}) &= \text{RollPolyHash}(g_{i,n}) \% |E_{n}^{hash}|
\end{align}

We adjust $e_i$ by dividing by the sum of the number of $n$-gram sizes and one, and employ RollPolyHash as specified in Appendix~\ref{appendix:rollpolyhash}. Section~\ref{section:ablations} investigates the impact of varying both the values of $n$ and the size of the embedding table in $n$-gram hash embeddings on flop-controlled scaling law behavior. Beyond hashing-based $n$-gram embeddings, our experiments also included frequency-based $n$-gram embeddings, with a comprehensive discussion provided in Appendix~\ref{appendix:ngrams}.

\subsubsection{Encoder Multi-Headed Cross-Attention}
\label{sec:enc-cross-attn}
Our approach largely adopts the cross-attention mechanism used in the input module of the Perceiver architecture~\citep{jaegle2021perceiver}, with a notable distinction: instead of employing a predetermined set of latent representations, our method utilizes patch-specific latent representations~(\autoref{fig:crossattn}), restricting attention solely to the bytes contained within each respective patch. In this module, each patch $p_j$ is associated with a query vector, which is produced by pooling the byte-level embeddings of $p_j$ and then applying a linear transformation, $\mathcal{E}_{C} \in \mathbb{R}^{h_{\mathcal{E}} \times (h_{\mathcal{E}} \times U_{\mathcal{E}})}$, where $U_{\mathcal{E}}$ represents the number of heads in the encoder’s cross-attention mechanism. Specifically, defining $f_{\text{bytes}}(p_j)$ as the ordered byte sequence corresponding to patch $p_j$, the computation proceeds as follows

\begin{align}
P_{0,j} &= \mathcal{E}_{C}(f_\text{bytes}((p_j)), f~ \text{is a pooling function} \\
P_l &= P_{l-1} + W_o\left(\text{softmax}\left(\frac{QK^T}{\sqrt{d_k}}\right)V\right) \\
\text{where } Q_j &= W_q(P_{l-1,j}), K_i = W_k(h_{l-1,i}), V_i = W_v(h_{l-1,i}) \\
h_l &= \text{Encoder-Transformer-Layer}_l(h_{l-1})
\end{align}

Here, $P \in \mathbb{R}^{n_p \times h_{\mathcal{G}}}$ denotes the set of $n_p$ patch embeddings that are intended for input to the global model. These embeddings are constructed by aggregating the byte-level embeddings $e_i$ corresponding to each patch $p_j$. The matrices $W_q$, $W_k$, $W_v$, and $W_o$ serve as linear projections for the queries, keys, values, and output, respectively. Specifically, the keys and values are generated by projecting the byte-level hidden states $h_i$ produced by the previous layer, or by $e_i$ in the case of the first layer. A patch-wise masking approach is adopted so that each query $Q_j$ is restricted to attend only to keys and values from the bytes within its respective patch $j$. Given that the model employs multi-head attention over $Q, K,$ and $V$, and since patch embeddings typically have a higher dimensionality ($h_{\mathcal{G}}$) compared to $h_{\mathcal{E}}$, during cross-attention $P_l$ is represented as multiple attention heads each of size $h_{\mathcal{E}}$. These are subsequently concatenated to form embeddings with dimensionality $h_{\mathcal{G}}$. A pre-LayerNorm operation is applied to queries, keys, and values before attention computation, and positional encodings are omitted from this cross-attention layer. A residual connection is also applied around the cross-attention component.

\subsection{Local Decoder} 

Analogous to the local encoder, the local decoder $\mathcal{D}$ utilizes a lightweight transformer architecture, comprising $l_{\mathcal{D}} << l_{\mathcal{G}}$ layers. This model reconstructs the original byte sequence $y_i$ from the sequence of global patch embeddings $o_j$. During this process, the local decoder generates raw bytes sequentially, conditioning each prediction on the sequence of previously decoded bytes. To facilitate this, it receives as input the hidden states generated by the local encoder for the corresponding byte sequence. The decoder alternates between $l_{\mathcal{D}}$ layers of cross-attention and transformer operations. Specifically, cross-attention precedes each transformer block, allowing the decoder to derive byte-level representations from the higher-level patch representations; these byte representations are then processed by the subsequent transformer layers within the local decoder.

\subsubsection{Decoder Multi-headed Cross-Attention} 

Within the decoder's cross-attention mechanism, the positions of queries and key/values are swapped: here, byte-representations serve as queries, while patch representations act as the key/values. For initializing the byte-representations in cross-attention, the model uses byte embeddings derived from the encoder’s final layer, specifically $h_{l_{\mathcal{E}}}$. For each subsequent decoder layer $l$, the byte-representations, $d_{l,i}$, are computed in the following manner:

\begin{align}
D_0 &= h_{l_{\mathcal{E}}} \\
B_l &= D_{l-1} + W_o\left(\text{softmax}\left(\frac{QK^T}{\sqrt{d_k}}\right)V\right), \\
  &\text{where } Q_i = W_q(d_{l-1,i}), K_i = W_k(\mathcal{D}_C(o_j)), V_i = W_v(\mathcal{D}_C(o_j)) \\
  D_l &= \text{Decoder-Transformer-layer}_l(B_l)
\end{align}

Here, $W_k$ and $W_v$ denote the key and value projection matrices, respectively. These matrices are utilized in performing a linear transformation followed by a split operation $\mathcal{D}_C$, both of which are applied to the final patch representations $o_j$ generated by the global model. The matrix $W_q$ serves as the query projection matrix and is applied to byte-level representations $d_{l-1}$ from the previous layer of the decoder transformer (or to $h_{l_{\mathcal{E}}}$ when considering the initial layer). The output of the cross-attention operation is projected through $W_o$, the output projection matrix, producing $B \in \mathbb{R}^{h_{\mathcal{D}} \times n_b}$, with $n_b$ representing the total number of output bytes. Subsequently, the next set of decoder representations, $D_l$, are derived by applying a decoder transformer layer to $B$, the result of the cross-attention block. Analogous to the cross-attention employed in the local encoder, we implement multi-head attention, incorporate pre-LayerNorm, exclude positional embeddings, and add a residual connection around the cross-attention component.

\section{Experimental Setup}
\label{section:experiments}
Our experimental design consists of rigorously controlled comparisons between {\textsc BLT} and models based on tokenization, ensuring that {\textsc BLT} is not afforded any potential benefits stemming from the use of extended sequence contexts.

\subsection{Pre-training Datasets} In our experiments, all model sizes are pre-trained on two primary datasets: 1) The Llama 2 dataset~\citep{touvron2023llama}, encompassing 2 trillion tokens amassed from an assortment of publicly accessible sources that have undergone rigorous cleaning and filtering to enhance data quality; and 2) {\textsc BLT}-1T: a newly assembled dataset containing 1 trillion tokens, curated from diverse public resources and incorporating a portion of the pre-training dataset made available by Datacomp-LM~\citep{li2024datacomp}. The Llama 2 dataset serves as the foundation for scaling law studies—specifically, to investigate the optimal token count in line with findings from~\cite{dubey2024llama} and to inform architectural decisions for {\textsc BLT}. In contrast, {\textsc BLT}-1T is utilized for full-scale pre-training, enabling direct comparisons with Llama 3 across various downstream benchmarks. Importantly, neither dataset contains data sourced from Meta’s products or services. Additionally, for the baseline tokenizer-based models, our experiments utilize the Llama 3 tokenizer, which offers a vocabulary of 128K tokens and consistently yielded improved baseline results compared to the Llama 2 tokenizer.

\subsection{Entropy Model}  
For our experiments, the entropy model comprises a byte-level language model that is trained on the identical data distribution as the complete {\textsc BLT} model. By default, this model is implemented as a transformer consisting of 100 million parameters, 14 transformer layers, a hidden size of 512, and employs a sliding window attention span covering 512 bytes. All other hyperparameter settings mirror those utilized in both the local and global transformer variants. As described in \autoref{section:ablations}, we evaluated the impact of altering model scale, architecture, and receptive field. Notably, if the model's receptive field is sufficiently constrained, the learned entropy model can be represented compactly through an efficient lookup table.

\subsection{Entropy Threshold and Equalizing Context Length} In models employing entropy-based patching, we determine a threshold for patching that results in a target average \textit{patch size} across the pretraining data mix. In {\textsc BLT}, the \textit{patch size} is a flexible parameter, unlike tokenization, and its selection greatly affects the effective context size accessible to the model. To control for these differences and to prevent models with larger patch sizes from having an unwarranted benefit due to longer contexts, we standardize the expected number of bytes processed in each batch. Consequently, for models with increased patch sizes, we compensate by decreasing the sequence length. Specifically, when training on Llama 2 data, a context window of 8k bytes is adopted; for the {\textsc BLT}-1T dataset, we expand the context to an average of 16k bytes, but the batch size remains constant at approximately 16M bytes throughout.

Although the mean batch size remains unchanged, dynamic patching strategies result in variable ratios of bytes to patches during batch loading. To enhance efficiency, our {\textsc BLT} training setup organizes batches by grouping patches, thereby minimizing the need for padding operations in the computationally demanding latent transformer. This approach maintains a fixed number of patches in each batch. For the Llama 2 and {\textsc BLT}-1T datasets, byte sequences are either padded or truncated to 12k and 24k bytes respectively throughout training, preventing excessive memory usage due to exceptionally large patch sequences.

\subsection{Entropy Model Context}
\label{section:entropy-context}
Our experimental observations indicate that applying entropy patching leads to incrementally larger patches, particularly in highly structured content such as multiple choice questions (as demonstrated in the MMLU example in \autoref{fig:mmlu_patching}), where repetition is common. The root of this variation lies in the decreased entropy of the recurring elements within the entropy model context. Therefore, for the comprehensive application of {\textsc BLT}-Entropy with a patch size of 4.5, we refresh the entropy context by inserting new lines and employ an approximate monotonicity constraint, since this strategy is less susceptible to "entropy drift" that may result from fluctuations in context length. Notably, this modification is restricted to entropy computation; our method for determining the entropy threshold remains unchanged.

\subsection{FLOPs Estimation} 
\label{section:flops}

Our approach closely adheres to the methodology established in Chinchilla~\citep{hoffmann2022training} for calculating transformer flops, which includes contributions from the feed-forward modules, qkvo projections within the self-attention mechanism, as well as the evaluation of attention and the final output projection. However, we depart from this by considering the input embedding layer as a highly efficient lookup operation rather than as a dense matrix multiplication; consequently, this component is treated as incurring zero flops. Consistent with prior studies, we assume that the backward pass requires twice the number of flops as the forward computation.

To determine the flops \textit{per byte} for {\textsc BLT} models, we sum the floating point operations for each component: the local encoder transformer, the global latent transformer, the local decoder transformer, as well as the cross-attention modules present in both encoder and decoder. This is expressed as:

\begin{align}
    \text{FL}_{\text{{\textsc BLT}}} &= \text{Transf. FL}(h_{\mathcal{G}}, l_{\mathcal{G}}, m=n_{ctx}/n_p,V=0) / n_p\\
   &+\text{Transf. FL}(h_{\mathcal{E}}, l_{\mathcal{E}}, m=w_{\mathcal{E}}, V=0) \\
   &+ \text{Transf. FL}(h_{\mathcal{D}}, l_{\mathcal{D}}, m=w_{\mathcal{D}}, V=256) \\ 
    &+ \text{Cross Attn. FL}(h_{\mathcal{E}}, l_{\mathcal{E}}, m=n_p, r=n_p/k) \times k/n_p \\ 
   &+ \text{Cross Attn. FL}(h_{\mathcal{D}}, l_{\mathcal{D}}, m=k, r=k/n_p) 
\end{align}

In this context, $n_{ctx}$ represents the sequence length measured in bytes, $n_p$ denotes the patch size, $r$ refers to the proportion of queries relative to keys and values, and $k$ specifies the ratio between the patch dimension and the byte dimension, which equates to the number of local model partitions combined to create a unified global representation (with $k=2$ as illustrated in \autoref{fig:crossattn}). The term $V$ indicates the size of the vocabulary utilized in the output projection, which appears solely within the local decoder. The attention module utilizes a context length $m$, which varies according to whether it is operating on a byte sequence or a patch sequence. Accordingly, we adjust the computation of attention FLOPs for each respective component. Detailed formulations for calculating FLOPs in Transformer-FLOPs and Cross-Attention FLOPs are provided in Appendix~\ref{section:flops-appendix}.

\subsection{Bits-Per-Byte Estimation} Perplexity serves as a token-level uncertainty metric and thus is inherently tied to a given tokenizer. To facilitate comparisons across both byte-level and token-level models, we adopt the Bits-Per-Byte (BPB) metric, as established in prior literature~\citep{xue2022byt5, yu2023megabyte, wang2024mambabyte}. BPB offers a tokenizer-agnostic alternative to perplexity. The metric is calculated as follows:

\begin{align}
    \text{BPB}(x)&= \frac{\mathcal{L}_{CE}(\pmb{x})}{\ln(2)\cdot n_{\text{bytes}}}
\end{align}

Here, the aggregate cross-entropy loss over the sequence $\pmb{x}$ quantifies the model's uncertainty, which is then divided by both the number of bytes in $\pmb{x}$ and a constant factor to yield a normalized score.

\subsection{Transformer Architecture Hyperparameters} For the transformer blocks utilized in {\textsc BLT}—encompassing both local and global models—we adopt an architectural framework consistent with Llama~3~\citep{dubey2024llama}. Specifically, feed-forward networks are equipped with the SwiGLU activation function~\citep{swiglu}, and self-attention mechanisms incorporate rotary positional embeddings (RoPE)~\citep{RoPe} with $\theta=500000$~\citep{xiong2024effective}. RMSNorm \citep{RMSNorm} serves as the layer normalization method throughout. In all self-attention layers where attention masks are static, such as those employing \textit{block causal} or \textit{fixed-window block causal} masking, we apply Flash attention~\citep{dao2022flashattention}, setting the window size to 512 for fixed-width attention patterns. For cross-attention operations that require dynamic patch-based masks, we make use of Flex Attention\footnote{\url{https://pytorch.org/blog/flexattention}}, enabling fused computation and substantial training speed improvements.

\subsection{{\textsc BLT}-Specific Hyperparameters} To evaluate the performance of {\textsc BLT} models, our experimental setup involves two primary approaches: examining scaling properties and assessing downstream tasks. We analyze models at several parameter scales—specifically, 400M, 1B, 2B, 4B, and 8B. Details regarding model architecture hyperparameters are listed in Appendix~Table~\ref{tab:arch}. For query initialization in the first cross-attention layer within the local encoder, we employ max-pooling. Each {\textsc BLT} model is configured to use $500,000$ hashes generated with a single hash function and n-gram lengths spanning from 3 to 8. The learning rate across all models is set to $4e-4$. This value is selected based on a hyperparameter sweep in the $1e-3$ to $1e-4$ range at the 400M and 1B scales, which revealed that an identical learning rate is optimal for both token and {\textsc BLT} models. For scaling experiments using Llama-2 data, batch sizes are chosen following the guidance in~\cite{dubey2024llama} or matched equivalently by byte count. The AdamW optimizer~\citep{adamw} is used, with parameters $\beta_1 = 0.9$, $\beta_2 = 0.95$, and $\epsilon=10^{-8}$. Training includes a linear warm-up of 2000 steps followed by a cosine learning rate decay schedule toward zero. Additionally, a weight decay of 0.1 is applied and global gradient clipping is enforced at a maximum norm of 1.0.

\section{Scaling Trends}
\label{section:scaling}

We provide a comprehensive analysis of the scaling behaviors exhibited by byte-level models, with the goal of guiding the further scaling of {\textsc BLT} models. Our scaling investigation seeks to overcome the shortcomings of earlier studies on byte-level models by: (a) examining scaling dynamics under compute-optimal training conditions, (b) training equivalent 8B parameter models on substantial datasets (up to 1T tokens or 4T bytes) and assessing their performance on downstream tasks, and (c) evaluating scaling characteristics in environments where inference costs are held constant. In a subsequent section, we delve into particular benefits associated with modeling byte-sequences.

\subsection{Parameter Matched Compute Optimal Scaling Trends}
\label{section:parameter-matched-scaling}
Leveraging the Llama 2 dataset, we train several \textit{compute-optimal} bpe and {\textsc BLT} models at four distinct model scales, spanning from 1B to 8B parameters. We visualize the relationship between the number of training flops and language modeling accuracy on a representative portion of the data mixture. The bpe models employ the parameter-to-data ratio established as optimal by Llama~3~\citep{dubey2024llama}, ensuring that, under a fixed training compute constraint, the models reach peak performance on the training corpus~\citep{hoffmann2022training}. This forms a solid reference point for comparison. In parallel, for each bpe model, we construct a {\textsc BLT} model trained on identical data, utilizing a Latent Transformer whose architecture and size are aligned with the corresponding bpe Transformer.

As shown in~\autoref{fig:scaling-compute-opt} (right), {\textsc BLT} models perform on par with or surpass their bpe equivalents, and this pattern persists across increasing model sizes and computational budgets. To our knowledge, {\textsc BLT} represents the first byte-level Transformer architecture to demonstrate scaling behavior comparable to BPE-based models in compute-optimal settings. These findings support our hypothesis that the optimal parameters-to-compute ratio observed for bpe models also extends to {\textsc BLT}, or is at least closely aligned.

Both advancements in architecture and the application of dynamic patching are essential to achieve scaling trends comparable to those observed with bpe. In ~\autoref{fig:scaling-compute-opt} (left), we present a comparison between models utilizing space-patching methods and Llama 3. To approximate SpaceByte \citep{slagle2024spacebyte}, we implement the {\textsc BLT} space-patching framework, omitting n-gram embeddings and cross-attention mechanisms. While SpaceByte demonstrates gains relative to Megabyte, it still falls significantly short of matching Llama 3’s performance. The right side of \autoref{fig:scaling-compute-opt} highlights the combined benefits arising from architectural modifications and the introduction of dynamic patching. At scale, {\textsc BLT} models achieve performance comparable to leading tokenizer-based architectures, including Llama 3.

Additionally, we note that the specific tokenizer selected impacts the results for tokenizer-based systems; models trained with the Llama-3 tokenizer consistently surpass those trained with the Llama-2 tokenizer under identical data conditions.

In summary, the {\textsc BLT} architecture occupies an intermediate position between Llama 2 and 3 in terms of performance when employing much larger patch sizes. While the average token sizes for Llama 2 and 3’s bpe tokenizers are 3.7 and 4.4 bytes, respectively, {\textsc BLT} demonstrates comparable scaling characteristics with patch sizes averaging 6 or even 8 bytes. Since inference flop requirements are inversely related to average patch size, adopting an 8-byte patch size yields approximately a 50\% reduction in inference flops. Moreover, models configured with larger patch sizes tend to exhibit enhanced performance as both model and dataset sizes increase. For example, {\textsc BLT} with an 8-byte patch size initially performs much worse than bpe Llama 2 at the 1B scale, but surpasses the bpe counterpart by the time it reaches the 7B scale. This pattern implies that larger patch sizes might yield even greater advantages at increased model scales, and raises the possibility that even bigger patch sizes could become practical as model capacity and training computation continue to expand.

\subsection{Beyond Compute Optimal Task Evaluations}
\label{section:evals}
To further investigate scaling dynamics, we train an 8B {\textsc BLT} model that surpasses the compute optimal scaling law using the {\textsc BLT}-1T dataset—a larger, higher-caliber corpus—and assess its capabilities across a diverse set of established classification and generation benchmarks. The evaluation suite comprises tasks targeting common sense reasoning, factual knowledge, and code generation:
\paragraph{Classification tasks} consist of ARC-Easy (0-shot)~\citep{Eval_ARC}, ARC-Challenge (0-shot)~\citep{Eval_ARC}, HellaSwag (0-shot)~\citep{Eval_hellaswag}, PIQA (0-shot)~\citep{Eval_piqa}, and MMLU (5-shot)~\citep{hendrycks2020measuring}. For these, we utilize a prompt-based likelihood scoring approach over the response options, reporting mean accuracy across each task.
\paragraph{Coding related generation tasks:} We evaluate code generation performance using pass@1 metrics on MBPP (3-shot)~\citep{Eval_MBPP} and HumanEval (0-shot)~\citep{Eval_HumanEval}, reflecting model proficiency in producing correct Python code completions.

In \autoref{tab:evals}, we present a comparison among three models, all trained using the {\textsc BLT}-1T dataset: a model leveraging the Llama 3 tokenizer with bpe encoding,\footnote{We selected the Llama 3 tokenizer with a 128k vocabulary size due to its superior performance relative to the 32k vocabulary in Llama 2.} as well as two distinct {\textsc BLT} variants. The first variant applies a space-patching approach ({\textsc BLT}-Space), while the second adopts an entropy-driven patching strategy ({\textsc BLT}-Entropy), incorporating an approximate monotonicity constraint and resetting the entropy model’s context with new line characters, as elaborated in~\autoref{section:entropy-context}. Each of these models is allocated an identical flop budget for training. Notably, for the {\textsc BLT}-Entropy model, we lower the entropy threshold from 0.6 to 0.1 during inference to enhance task accuracy, though this leads to a greater number of inference steps.

The {\textsc BLT}-Entropy model demonstrates superior performance compared to the Llama 3 model in 4 of the 7 evaluated tasks, despite both being trained on an equal amount of data (in bytes). This enhanced performance can likely be attributed to two factors: (1) the more efficient utilization of training computation made possible by dynamic patching, and (2) the explicit modeling of information at the byte level rather than relying on token-based representations.

Conversely, {\textsc BLT}-Space exhibits lower performance than the Llama 3 tokenizer in all tasks except one; however, its average patch size of 6 bytes results in substantially fewer inference flops. By contrast, both the bpe and entropy-patching based models demonstrate similar average patch sizes, each around 4.5 bytes on the combined training dataset. Given an identical training budget, the model employing the larger patch size is able to process 30\% more data than the other models. This increased data coverage may further distance {\textsc BLT} from achieving compute-optimal efficiency.

\subsection{Patches Scale Better Than Tokens} Utilizing {\textsc BLT} models enables concurrent scaling of both model and patch sizes, all while preserving a consistent computational (flop) cost for training and inference, and without changing the quantity of training data. Unlike fixed-vocabulary token-based architectures, patch-based models provide the flexibility to expand patch size arbitrarily, a property that liberates them from the inherent efficiency compromises mentioned in Section~\ref{section:bpe-patch}. Using larger patches reduces computational overhead, which in turn allows for the expansion of the global latent transformer, as it is invoked less frequently.

We perform a fixed inference scaling analysis to evaluate whether increasing model size while reducing step count on larger patches yields superior performance compared to smaller models executing more steps. Using initial model sizes of 400 million and 3.6 billion parameters based on the Llama 2 tokenizer, we identify flop-matched models incorporating the Llama 3 tokenizer, as well as {\textsc BLT}-Entropy architectures with average patch sizes of 6 and 8 bytes, all trained on the same datamix (model specifications are listed in~\autoref{tab:fixed-inf-params}). Notably, for models with a patch size of 8, we employ 3 encoder layers rather than 1. Each configuration is trained across a range of training flop budgets.

\autoref{fig:fixed_inference_scaling} demonstrates that {\textsc BLT} architectures exhibit more favorable scaling behaviors than tokenization-based models across both categories of inference flops. While BPE-based models demonstrate superior performance at lower training budgets, they are soon outperformed by {\textsc BLT} models as training expenditures move modestly beyond the compute-optimal threshold. In real-world applications, it may be advantageous to allocate greater resources for the one-time pretraining phase to obtain higher-performing models under a fixed inference cost. The 8B parameter class—exemplified by models such as Llama 3.1, which is pretrained on data volumes two orders of magnitude greater than the compute-optimal amount for its size—illustrates this approach.

The point at which {\textsc BLT} outperforms token-based models shifts somewhat nearer to the compute-optimal point as models in the higher flop class are considered (changing from being 3x to 2.5x over the compute-optimal allocation). Likewise, the model using patch size 8 in the larger flop class demonstrates a steeper scaling curve and surpasses the other models at an earlier stage. As outlined in Section~\ref{section:parameter-matched-scaling}, larger patch sizes tend to show performance more comparable to BPE models at increased scales. This trend may be partially due to the fact that, as scale increases, the relative proportion of total flops consumed by the byte-level Encoder and Decoder modules—which tend to scale less rapidly than the Latent Transformer—diminishes. For instance, scaling total parameters twenty-fold from 400M up to 8B only leads to an approximate doubling of the local parameters within {\textsc BLT}. This relationship is significant because the impact of larger patch sizes is confined to the flops from the patch Latent Transformer and does not affect the byte-level components. Indeed, this explains why {\textsc BLT}-Entropy with ps=8 changed from requiring 1.6x to 1.7x the size of the Llama 2 model as the model scale was increased.

In conclusion, our investigation into patch-length scaling reveals that the {\textsc BLT} patch-based architecture benefits from enhanced scaling behavior when both patch size and overall model size are increased together. These favorable scaling patterns not only hold but may become more pronounced as the model scale grows.

\section{Byte Modeling Improves Robustness} We further evaluate the robustness of {\textsc BLT} in contrast to token-based models that do not natively process byte-level data, and introduce a technique for converting pretrained token-based models to operate at the byte level.
\label{sec:robustness}

\subsection{Character-Level Tasks}

One of the initial drivers for developing byte-level models was their inherent resilience to noise present at the byte level in input data, as well as their ability to recognize sub-token components—an area where tokenizer-based approaches often face difficulties. To systematically assess these capabilities, we conduct supplementary experiments using benchmarks specifically designed to test resistance to noisy inputs and sensitivity to character-level details, encompassing English, a variety of languages, numerals, and phonetic symbols. The outcomes of these experiments are summarized in Table~\ref{tab:char_tasks}.

\paragraph{Noisy Data} To assess the robustness of tokenizer-based models versus {\textsc BLT}, we generate noisy variants of the benchmark classification tasks detailed in Section~\ref{section:evals}. Five character-level noise techniques are employed to introduce perturbations: (a) \textit{AntSpeak}: The text is rendered in all uppercase letters, separated by spaces between each character. (b) \textit{Drop}: 10\% of the text's characters are removed at random. (c) \textit{RandomCase}: Half of the text’s characters are randomly set to uppercase while the remaining half are set to lowercase. (d) \textit{Repeat}: 20\% of characters are repeated up to four times. (e) \textit{UpperCase}: Every character in the text is converted to uppercase. For evaluation, we apply each noise approach to the prompt, the completion, or both independently, and average the results over these scenarios. Table~\ref{tab:char_tasks} displays results on noised HellaSwag~\citep{Eval_hellaswag}, revealing that {\textsc BLT} consistently demonstrates superior robustness compared to tokenizer-based approaches—showing an average improvement of 8 points relative to the model trained on the same dataset and even outperforming the Llama 3.1 model, which was trained with considerably more data.

\paragraph{Phonology - Grapheme-to-Phoneme (G2P)} We evaluate how effectively {\textsc BLT} converts a sequence of graphemes (the characters that compose a word) into its corresponding phonemic transcription. Table \ref{tab:char_tasks} displays performance on the G2P task under a 5-shot scenario, utilizing Phonology Bench~\citep{suvarna2024phonologybench}. Our analysis shows that {\textsc BLT} achieves superior results compared to the Llama 3 1T model that uses a tokenizer-based approach.

\paragraph{CUTE}
To evaluate the character-level capabilities of {\textsc BLT}, we employ the CUTE benchmark~\citep{edman2024cute}, which encompasses multiple tasks organized into three main areas: composition comprehension, orthographic similarity recognition, and sequence manipulation skills. This suite of tasks is notably challenging for most models that utilize token-based tokenization, as such models generally encode information about token spellings but have difficulty applying this knowledge to perform manipulations at the character level. As shown in Table~\ref{tab:char_tasks}, {\textsc BLT}-Entropy surpasses both BPE-based Llama 3 variants by over 25 points in overall CUTE performance. Notably, our model exhibits near-perfect results in character manipulation, achieving 99.9\% accuracy in both spelling assessments. These substantial gains are achieved despite {\textsc BLT} being exposed to only one-sixteenth the training data used for Llama 3.1, highlighting the challenges BPE tokenizers face in learning character-level representations. Figure~\ref{fig:cute_examples} presents selected examples where Llama 3's tokenizer-based approach falls short, whereas {\textsc BLT} excels. The only areas where BPE tokenizers outperform are word deletion and insertion, which likely present inherent difficulties for byte-level models; however, the margin remains small. Transitioning from characters to words may ultimately be less complex than the reverse. Our evaluation employs a consistent protocol across all tasks, utilizing the original prompts from Huggingface. It is possible that prompt engineering could further improve BPE models' performance in these benchmarks.

\paragraph{Low Resource Machine Translation} We assess the performance of {\textsc BLT} in both directions of translation—into and out of English—across twenty-one lower-resource languages representing six widely spoken language families, each with different writing systems. This evaluation is based on the FLORES-101 benchmark~\citep{goyal2022flores}, and the corresponding SentencePiece BLEU scores are provided in Table~\ref{tab:flores}. The empirical evidence indicates that {\textsc BLT} surpasses a baseline model utilizing the Llama 3 tokenizer, obtaining a 2-point overall improvement when translating into English, and a 0.5-point gain when translating from English. In language pairs with substantial available resources, {\textsc BLT} delivers results similar to, or marginally exceeding, those of Llama 3. More notably, {\textsc BLT} achieves consistently stronger outcomes over Llama 3 for a variety of language pairs within less-resourced language families, highlighting the strengths of byte-level modeling in handling diverse and uncommon byte sequences.

\subsection{Training {\textsc BLT} from Llama 3}
\label{sec:distillation}
We investigate a training strategy in which {\textsc BLT} models benefit from pre-existing pretrained tokenizer-based models, thus facilitating more efficient and rapid convergence. Specifically, we initialize the global transformer parameters in {\textsc BLT} with those of a pre-trained Llama 3.1. During subsequent training, the global transformer's weights are fine-tuned using a learning rate set at one-tenth of that assigned to the local encoder and decoder components, and training proceeds for the same optimal number of steps as Llama 3. The performance of this approach is benchmarked against a baseline {\textsc BLT} model in Table~\ref{tab:distill}. The results demonstrate that {\textsc BLT} initialized from Llama 3.1 delivers markedly better results than both the standard Llama 3 and baseline {\textsc BLT} models, all under identical flop budgets. Furthermore, when contrasted with the {\textsc BLT}-Entropy variant described in Table~\ref{tab:evals}, which was trained on a much more extensive dataset (1T tokens), the {\textsc BLT} derived from Llama 3.1 continues to outperform on the MMLU task. This finding indicates that this transfer approach is highly effective in minimizing required training computation.

This approach can alternatively be interpreted as converting models that rely on tokenizers into models that operate without them, thereby transforming a pre-trained LLaMA 3.1 model into a {\textsc BLT} variant. For a thorough evaluation, Table \ref{tab:distill} presents the performance of the original LLaMA 3.1, which has been trained on 15T tokens, and compares it to that of the {\textsc BLT} variant derived from LLaMA 3. While our adapted model shows only a minor decrease in performance on MMLU and HumanEval, there is a notably larger reduction in effectiveness on other benchmarks. These outcomes indicate that additional research is necessary to better exploit the capabilities of the pre-trained model, with particular attention to refining data composition and tuning relevant hyperparameters.

\section{Ablations and Discussion}
\label{section:ablations}
This section presents ablation studies that support the architectural decisions behind {\textsc BLT}, as well as details regarding the patching approach and the selection of hyper-parameters for the 8B parameter {\textsc BLT} model trained using the {\textsc BLT}-1T dataset.

\paragraph{Entropy Model Hyper-parameters} To analyze the impact of entropy model size and context window length on model scaling behavior, we trained several byte-level entropy transformer variants, ranging from 1 million to 100 million parameters and employing context windows of 64 up to 512 tokens. Figure \autoref{fig:entropy_ablation} illustrates the scaling curves (bits-per-byte vs training FLOPs) generated with our $400m$ and $1b$ {\textsc BLT} models, which were trained using the Llama-2 dataset. Our experiments demonstrate that scaling both model size and context window leads to improved performance; however, the benefits become marginal once the model surpasses 50 million parameters.

\paragraph{Types of Patching}

We conduct an ablation study on the four patching methods described in Section~\ref{section:patching}, namely: 1) Strided Patching with strides of 4 and 6, 2) Whitespace-based Patching, 3) BPE Tokenizer patching utilizing the Llama 3 tokenizer, and 4) Entropy-driven patching informed by a small byte-level language model.

Although dynamic patching shortens the practical sequence lengths, we regulate the sequence length to ensure a consistent context window across different patching strategies. Each model is exposed to an equivalent expected number of bytes per sequence during both training and inference phases, thereby minimizing confounding effects due to unequal context modeling capacities. As shown in Figure~\ref{fig:scaling-compute-opt}, these ablation experiments reveal that all alternative patching strategies surpass static patching in performance, with space patching performing nearly as well as dynamic entropy-based patching.

In \autoref{tab:evals_ablation}, we show comparative benchmark results for {\textsc BLT} models using three different approaches: models based on tokenizers, models employing space patching, and models utilizing entropy-based patching, all trained on the Llama 2 dataset for a number of steps determined to be optimal~\citep{dubey2024llama}. While space patching offers a more straightforward method since it does not require invoking an entropy model dynamically during the training phase, our analysis indicates that the improvements demonstrated by entropy-based patching on scaling behaviors~(Section~\ref{section:scaling}) persist and remain evident on downstream benchmark evaluations as well.\footnote{Space patching results are from earlier experimental runs lacking cross-attention, but we observe analogous patterns even when cross-attention is included.}

\paragraph{Cross-Attention} 
In~\autoref{tab:crossattn_ablations_1b}, we conduct an ablation study evaluating the impact of incorporating cross-attention at different locations within the encoder and decoder of {\textsc BLT}. For encoder-side cross-attention, we experiment with initializing the queries in three ways: (1) assigning an identical learned embedding to each global state, (2) employing a hash embedding derived from the bytes contained in the patch, and (3) using a pooled representation from the encoder’s hidden states over the patch bytes at the corresponding encoder layer.

Our results indicate that incorporating cross-attention within the \textit{decoder} yields the highest performance gains. Within the encoder, the application of cross-attention offers minor benefits, which occur primarily when queries are initialized via pooling. Moreover, cross-attention proves especially advantageous on the Common-Crawl dataset, with the improvements being more pronounced when utilizing larger patch sizes.

\paragraph{n-gram Hash Embeddings} 

We experiment with n-gram hash embedding vocabulary sizes of 0, 100K, 200K, and 400K, and report the corresponding outcomes in Table~\ref{tab:ngram_ablations_1b}. Our observations show that the use of hash embeddings confers benefits across all evaluated domains, with a marked impact on Wikipedia and Github (where the bpb improvement reaches 0.04, in contrast to just 0.01 after 15k steps at 8B). When operating at the 8B scale, reducing hash vocabulary from 500K to 300K only marginally affected performance, changing by approximately 0.001 bpb after 15k steps. These findings suggest that hash embeddings are essential for enabling {\textsc BLT} to achieve parity with models based on tokenizers, though gains plateau once the number of hashes exceeds 300K. Furthermore, the data indicate that hash embedding improvements and those from cross-attention are largely additive, benefiting different datasets in distinct ways.

\paragraph{Local Model Hyperparamaters} Table \ref{tab:local_layers} presents an ablation study of different configurations for the number of layers in the local encoder and decoder. In combination with hash n-gram embeddings, {\textsc BLT} achieves strong performance when the encoder is kept very minimal—specifically, with only a single layer—while utilizing a more substantial decoder.

\section{Related Work}
\label{section:related_work}

\textbf{Character-Level RNNs:} The task of Character Language Modeling has long been of significant interest in the neural modeling community~\citep{sutskever2011generating,mikolov2012subword,graves2013generating}, primarily due to its natural ability to handle out-of-vocabulary terms without relying on traditional back-off strategies. In their work, \cite{kim2016character} develop a model that exclusively handles input as characters via convolutional and highway layers, which then serve as input to LSTM-based RNNs. This architecture achieves competitive results with the leading RNN-based language models of that era for English, while delivering superior performance in morphologically complex languages—a key benefit attributed to character-level large language models. Additionally, \cite{kenter2018byte} explore byte-level LSTM approaches for machine comprehension tasks, surpassing word-level models in performance, particularly for morphologically rich languages such as Turkish and Russian. Following this trend, \cite{zhang2015character} utilize character-based convolutional neural networks for classification problems, in some cases exceeding the efficacy of word-based models. Furthermore, \citet{chung2019hierarchical} propose a hierarchical LSTM framework, equipped with boundary-detection modules at each level, which facilitates the identification of hidden structural hierarchies in text, resulting in enhancements to character-level language modeling. In a related direction, ByteNet, introduced by \cite{kalchbrenner2016neural}, employs convolutional neural network layers on character inputs instead of leveraging attention mechanisms for tasks such as machine translation.

\textbf{Character-Level Transformers:} Advances in transformer architectures employing attention mechanisms~\citep{vaswani2017attention} alongside subword tokenization methods~\citep{sennrich-etal-2016-neural} have led to substantial gains in neural language modeling and related evaluation tasks. Nevertheless, relying on word or subword units inherently imposes a particular inductive bias, dictating the abstraction level at which the model functions. In an effort to reconcile the success of transformer frameworks with the early encouraging outcomes of character-level modeling, \citet{al2019character} employed extremely deep transformer models augmented with auxiliary loss functions. This approach enabled them to surpass previous LSTM-based character language models. Despite these improvements, there remained a considerable performance disparity when compared to language models operating at the word level. Similarly, GPT-2~\citep{radfordgpt2} found that, for extensive datasets such as the 1 Billion Word Benchmark, byte-level language models lagged behind their word-level counterparts in terms of competitiveness.

Although \cite{Choe2019BridgingTG} showed that transformer-based byte-level LLMs can surpass the performance of subword-level LLMs with similar parameter sizes, these models require substantially greater computational resources and extended training times. In a related vein, \cite{el2020characterbert} introduced a BERT variant called CharFormer that constructs word representations via convolutions over character embeddings; while this model yields better results in the medical domain, it does so at the expense of significantly increased computation. The work of \cite{clark2022canine} presents CANINE, an encoder-only model with 150M parameters that operates directly on character sequences. CANINE features a deep transformer stack analogous to the core structure of our global model, and integrates both a local transformer and strided convolutions for character input downsampling. This architecture achieves superior downstream performance compared to mBERT, the equivalent token-level encoder-only model, on a variety of multilingual tasks. ByT5~\citep{xue2022byt5} investigated byte-level encoder-decoder models that entirely forego patching operations. Although ByT5 exhibits increased robustness to noise and remains competitive with tokenizer-based counterparts using only one-fourth as much data, its lack of input patching requires expensive attention computations over each individual byte, which significantly escalates compute requirements. Switching from subword to byte-level modeling lengthens the input sequences, posing efficiency and scaling challenges for these models. More recently, leveraging the Mamba Architecture~\citep{gu2023mamba}, which can keep a fixed-size memory across extensive context windows, \cite{wang2024mambabyte} trained a byte-level Mamba model—again without patching—and demonstrated that it can exceed the performance of byte-level transformers in flop-controlled experiments at the 350M parameter scale, as measured by bits-per-byte on multiple datasets.

\textbf{Patching-based approaches:} Leveraging patching techniques has proven to be an effective means for reducing the excessive flop counts typical of byte-level LLMs, all while potentially maintaining their performance. Several studies have reported promising outcomes at relatively small scales in terms of both model size and training data. For instance, \cite{nawrot-etal-2022-hierarchical} conducted experiments with static downsampling and upsampling via patching, leading to the development of the hourglass transformer, which surpassed other byte-level baselines at the 150M parameter level. Building on this foundation, \cite{nawrot-etal-2023-efficient} introduced enhanced dynamic patching methods, incorporating an end-to-end trainable boundary-predictor, a variant supervised through specific tokenizers, and an entropy-informed patching technique resembling {\textsc BLT}. Their results demonstrated that such dynamic patching strategies outperformed standard transformer architectures in language modeling tasks with 40M parameters and roughly 400M training tokens. Additionally, \cite{lester2024training} explored the impact of training on sequences compressed with arithmetic coding to realize compression levels superior to those achievable by BPE. Using their equal-info windows method, they achieved better performance than byte-level baselines in language modeling tasks, though their approach did not exceed that of subword-based baselines.

This study builds upon and is most directly associated with MegaByte~\citep{yu2023megabyte}, a decoder-only causal LLM that employs a predetermined static patching mechanism and concatenates representations to transform bytes into patches. On the decoder end, it relies on a local model to convert these patches back to bytes. MegaByte has been shown to achieve performance comparable to tokenizer-based models at the 1B parameter level on datasets encompassing approximately 400B bytes. In all of our evaluations, we systematically compare against MegaByte and observe that its static patching approach falls short relative to the most recent, compute-efficiently trained tokenizer-based models under controlled FLOP budgets; we illustrate how {\textsc BLT} effectively closes this performance gap. \cite{slagle2024spacebyte} corroborate these findings, proposing an enhancement to the static patching strategy by incorporating patching over whitespaces and other space-type bytes and introducing a local encoder model. Their results demonstrate gains over tokenized transformer baselines in compute-constrained conditions for certain data sources, such as Github and arXiv, with 1B-parameter models. We include empirical results using this methodology and demonstrate that additional advances in architecture are necessary to further improve the scalability of byte-level models and to truly rival the latest tokenizer-based models, for example Llama 3.

\section{Limitations and Future Work}

In this study, we adhere to the optimal number of training steps established for Llama~3~\citep{dubey2024llama} when making architectural decisions. Nonetheless, since these scaling laws were derived based on BPE-level transformer models, they may not yield the most efficient (data, parameter size) proportions when applied to {\textsc BLT}. Determining scaling laws specifically suited to {\textsc BLT} remains an open question for subsequent research, and could potentially result in more advantageous scaling characteristics for our model. Furthermore, the majority of our experiments were performed at parameter counts of up to 1B, and it remains plausible that architectural optima might shift at greater scales—such as 8B parameters or more—potentially leading to superior results at these higher capacities.

Current transformer libraries and codebases are primarily optimized for transformer models that utilize tokenization. Although our work features flop-equivalent experimental comparisons and incorporates some optimized implementations (like FlexAttention) to manage components differing from the standard transformer design, the runtime performance of our models might still lag behind that of tokenizer-based approaches. Further improvements in efficiency could help bridge this gap in wall-clock execution time.

Whereas {\textsc BLT} relies on an entropy model for patching that is trained independently, an appealing avenue for future exploration is to develop a patching model that is jointly optimized in an end-to-end manner. In Section \ref{sec:distillation}, we offer preliminary empirical evidence that suggests it is feasible to "byte-ify" tokenizer-based architectures like Llama 3—models that have seen over 10T training tokens—by initializing the global transformer component with their pretrained weights and keeping these parameters fixed. Pursuing this line of research further may lead to strategies that not only preserve the strengths conferred by bytefying, but also improve upon the performance of existing tokenizer-based systems, all without necessitating their full retraining.

\section{Conclusion}
\label{section:conclusion}

In this work, we introduce the Byte Latent Transformer (\textbf{BLT}), a novel model architecture that challenges the standard reliance on fixed-vocabulary tokenization in large language models. Through the deployment of an adaptive, trainable approach to aggregating bytes into patches, {\textsc BLT} dynamically distributes computational effort according to the underlying data complexity. This mechanism yields marked gains in efficiency and resilience. Our large-scale empirical analysis indicates that {\textsc BLT} rivals the capabilities of leading tokenization-based systems such as Llama 3 when trained with up to 8B parameters and exposed to 4T bytes of data, while also being able to reduce inference flops by up to 50\% in exchange for slight compromises on evaluation benchmarks. Additionally, {\textsc BLT} paves the way for a new scaling strategy, where both model and patch sizes can be enlarged concurrently under constant inference constraints—a property especially relevant for real-world compute limitations. By operating directly on raw bytes, {\textsc BLT} enhances the handling of data with infrequent patterns, notably improving robustness against noisy input and deepening the model’s grasp of sub-word compositions. Collectively, our findings suggest that {\textsc BLT} offers a viable and compelling substitute to conventional tokenization-centric architectures, promoting a more flexible, scalable, and robust foundation for advanced language modeling.

\end{document}